\documentclass[10pt,a4paper]{amsart}
\usepackage[margin=19mm]{geometry}
\usepackage{amsmath,amssymb,mathtools}
\usepackage[scr=rsfs,cal=euler]{mathalfa}
\usepackage{xfrac}
\usepackage[unicode,hidelinks,bookmarks=false]{hyperref}
\newcommand{\ie}{i.e.\ }
\newcommand{\cf}{cf.\ }
\newcommand{\resp}{respectively\ }
\newcommand{\Subsection}{\subsection}
\providecommand{\nobreakdash}{\nobreak\hbox{-}}
\setlength{\emergencystretch}{2em}
\multlinegap0pt
\usepackage{braket}
\usepackage{amssymb}
\usepackage[shortlabels]{enumitem}
\usepackage[normalem]{ulem}
\newtheorem{lemma}{Lemma}
\title{Pulsation of quantum walk between two arbitrary graphs with weakly connected bridge}
\author{Taisuke Hosaka, Etsuo Segawa}
\date{Source: arXiv:2511.12872v2 (CC BY 4.0)}
\begin{document}
\maketitle

\noindent\emph{Source and licence.} Pulsation of quantum walk between two arbitrary graphs with weakly connected bridge. Version arXiv:2511.12872v2 (CC BY 4.0), distributed by arXiv under the Creative Commons Attribution 4.0 International licence, http://creativecommons.org/licenses/by/4.0/, as recorded in the arXiv licence metadata for this exact version and shown on the abstract page https://arxiv.org/abs/2511.12872v2. Full licence notice: \texttt{licenses/arxiv-2511.12872-CC-BY-4.0.txt}.\par\smallskip

\noindent\textit{Editorial scope.} Counted unit: the article's lemma overlap (source0.tex lines 629--637). The article's complete original proof of it is delivered with it (source0.tex lines 638--652, one \texttt{proof} environment of 15 lines). No mathematics is written, repaired or re-derived: the statement and the proof are the article's own lines, in the article's own notation, and no retained line is substituted or re-flowed.  Retained uncounted as the source-local context the proof needs: lines 182--185; lines 596--614.  Nothing was omitted from any retained span, so the package carries no omission record.  No retained line contains a citation, so the wrapper carries no bibliography and no bibliography entry.  No retained line contains a figure environment, an image-inclusion command or drawing code, so no image is delivered and the package declares no asset dependency.  Editorial formatting only, in addition: an AMS \texttt{amsart} preamble that restates the article's own declarations and macros used by the retained lines, namely the environments \texttt{lemma} on separate counters, together with the standard packages those lines need.  The retained environments print in this document as Lemma 1 in order of appearance, whereas the article prints the same items with the numbering of its own full document.  The complete native source of the article as distributed in the arXiv e-print for this exact version is bundled unchanged as \texttt{sources/189238/source0.tex}.
\begin{align}
    \label{eq: initial state}
    \ket{\psi_0}=\frac{1}{\sqrt{|A_1|}}\sum_{a\in A_1}\ket{a}.
\end{align}

\begin{lemma}
    \label{lem: eigenpair of U}
    Let $U(\epsilon)$ be the time evolution matrix. Then it follows that
    \begin{align*}
        \{1, e^{\pm i \theta(\epsilon)}\} \subset \mathrm{Spec}(U(\epsilon)).
    \end{align*}
    Corresponding eigenvectors $\ket{\psi_1} \in \mathrm{Ker}(I-U(\epsilon))$ and $\ket{\psi_{\pm\theta}} \in \mathrm{Ker}(e^{\pm i\theta(\epsilon)}-U(\epsilon))$ are given by
    \begin{align*}
        \psi_1(a)=\frac{1}{\sqrt{|A_1|+|A_2|}},
    \end{align*} 
    \begin{align*}
        \psi_{\pm\theta}(a)=\frac{1-e^{\pm i\theta(\epsilon)}}{\sqrt{2|A_1||A_2|(|A_1|+|A_2|)}\,|\sin (\theta(\epsilon))|}\times
        \begin{cases}
            -|A_2| &: t(a) \in V_1, \\
            |A_1| &: t(a) \in V_2,
        \end{cases}
    \end{align*}
    respectively.
\end{lemma}

\begin{lemma}
\label{lem: overlap}
Let $\ket{\psi_1}$ and $\ket{\psi_{\pm \theta}}$ be eigenvectors of $U(\epsilon)$ corresponding to eigenvalues 1 and $e^{\pm i\theta(\epsilon)}$, respectively.
Let $\ket{\psi_0}$ be an initial state given by Eq. (\ref{eq: initial state}). Then we have
    \begin{align*}
        |\braket{\psi_1|\psi_0}|&=\sqrt{\frac{|A_1|}{|A_1|+|A_2|}}, \\
        |\braket{\psi_{\pm\theta}|\psi_0}|&=\frac{1}{\sqrt{2}}\sqrt{\frac{|A_2|}{|A_1|+|A_2|}}+O(\epsilon).
    \end{align*}
\end{lemma}

\begin{proof}
    By directly, we get the first formula. Combining Eq.\,(\ref{eq: initial state}) with Lemma \ref{lem: eigenpair of U}, we see
    \begin{align*}
        |\braket{\psi_{\pm \theta}|\psi_0}|=\frac{1}{\sqrt{2}}\sqrt{\frac{|A_2|}{|A_1|+|A_2|}}\times \left| \frac{1-e^{i\theta(\epsilon)}}{\sin \theta(\epsilon)}\right|.
    \end{align*}
    Thus, we have
    \begin{align*}
        \left| \frac{1-e^{\pm i\theta(\epsilon)}}{\sin \theta(\epsilon)}\right|
        &= \left|\frac{e^{\pm i\theta(\epsilon)/2}(e^{\pm i \theta(\epsilon)/2}-e^{\mp i \theta(\epsilon)/2})}{\sin \theta(\epsilon)} \right| \\
        &= \left|\frac{2i \sin (\theta(\epsilon)/2)}{\sin \theta(\epsilon)}\right| \\
        &= \left|\frac{1}{\cos (\theta(\epsilon)/2)} \right| \\
        &= 1+O(\epsilon).
    \end{align*}
    Therefore, we get the desired conclusion.
\end{proof}

\end{document}
